\.section{Background}
\label{sec:rs-background}

The Recursive Spawning Protocol (RSP) operates at the intersection of several research and engineering traditions: agent lifecycle management in multi-agent AI systems, accountability and provenance in distributed computing, fixed versus mutable delegation structures, hierarchical task decomposition in AI, and the BX3 Framework's three-layer model as the architectural substrate for immutable parent chains. Each tradition contributes an essential conceptual component; RSP's contribution is integrating them into a unified structural solution to the accountability gap that emerges when agents spawn children autonomously and those children continue operating with real-world effects.

% -------------------------------------------------------
\subsection{Agent Lifecycle Management in Multi-Agent Systems}
% -------------------------------------------------------

Multi-agent systems (MAS) research has long recognized that an agent's lifecycle --- its creation, activation, operation, suspension, and termination --- is a first-class design concern, not merely an implementation detail. The foundational taxonomy introduced by Stone and Veloso classifies multi-agent systems along axes of agent heterogeneity, agent knowledge quality, environment accessibility, and the nature of agent interactions \cite{stone2000}. Within this taxonomy, recursive spawning is most relevant to cooperative, homogeneous systems operating in partially accessible environments --- precisely the regime in which modern agentic runtimes such as AutoGPT \cite{autoGPT2023}, LangChain \cite{langchain2023}, and Microsoft's AutoGen \cite{autogen2023} operate. Classical agent lifecycle management was primarily an engineering concern: spawning, scheduling, and terminating agents to maximize throughput or minimize latency, with accountability residing with the system operator who configured the agent population. As AI agents move into healthcare, infrastructure control, and financial decision-making, this assumption becomes untenable. When a spawned agent's action produces real-world consequences, the question of \emph{who authorized the spawn}, through what chain of delegation, and with what scope constraints is no longer academic --- it is a legal, regulatory, and safety question with no clean answer in conventional architectures.

Ferber and Müller formalized lifecycle interactions as a state machine coupled to the organizational context in which an agent operates, demonstrating that organizational roles --- not just individual agent states --- must be modeled to ensure system-wide coherence \cite{ferber1996}. Padgham and Thiel further established that incomplete lifecycle management, particularly the failure to distinguish between suspension and termination states, introduces subtle but systemic failures: orphaned or zombie agents that continue to consume resources and emit effects after their parent has exited \cite{padgham2001}. More recent work has surfaced the specific problem of unbounded spawning. Zhuge et al. document the agent chain problem: LLM-based agents that recursively spawn children to handle sub-tasks, without a defined termination condition or depth ceiling, can produce chains of dozens of intermediate agents whose aggregate authority differs substantially from what any single agent was authorized to exercise \cite{zhuge2024}.

Several recent protocols have begun addressing the accountability gap in multi-agent lifecycles directly. The Human Delegation Provenance (HDP) protocol identifies the core vulnerability in multi-hop delegations: in chains of the form human $\rightarrow$ orchestrator $\rightarrow$ sub-agent $\rightarrow$ tool-agent, verifying that terminal actions were genuinely authorized by a human requires reconstructing the full delegation chain --- a task that OAuth 2.0, JWTs, and UCAN do not adequately address for append-only, human-provenance requirements in agentic systems \cite{hdp-draft,hdp-draft-2}. HDP proposes cryptographic tokens capturing each delegation hop as a signed entry in an append-only chain, enabling offline verification using only the issuer's public key and session identifier. The Agent Lifecycle Protocol (ALP) extends this with canonical lifecycle events --- Genesis, Fork, Migration, Retirement, Succession, and Decommission --- with explicit state-machine semantics, transition rules, and a genealogical registry capturing both genetic (model, architecture) and epigenetic (configuration, memory, reputation history) lineage \cite{alp2025}. PROV-AGENT demonstrates convergence with the W3C PROV standard by extending it with the Model Context Protocol to link agent decisions and prompts with end-to-end workflow data, enabling near real-time provenance capture across facility boundaries \cite{prov-agent2025}. Unified Agent Lifecycle Management (UALM) provides a complementary five-layer blueprint for end-to-end governance of agent fleets: identity and persona registry, orchestration and cross-domain mediation, PHI-bounded context and memory, runtime policy enforcement with kill-switch triggers, and lifecycle management linked to credential revocation and audit logging \cite{ualm2025}. Together, these protocols establish that the industry has identified the accountability gap in agentic lifecycles and is converging on cryptographic, append-only delegation records as the operative solution.

% -------------------------------------------------------
\subsection{Accountability and Provenance in Distributed Systems}
% -------------------------------------------------------

The accountability gap in agentic systems maps directly onto a well-established problem in distributed systems: establishing trustworthy provenance for chains of computation and delegation when multiple actors contribute to a terminal outcome. The W3C PROV standard provides the foundational vocabulary --- entities, activities, and agents --- for modeling who did what, under whose authority, and with what effect \cite{prov-standard}. PROV's core insight is that provenance is a first-class structural concern, not a logging afterthought, and that the relationships between provenance elements encode accountability-relevant information.

In conventional distributed systems, accountability for delegated actions is handled through access control lists (ACLs), role-based access control (RBAC), and capability-based security models. These approaches share a common weakness in the agentic context: they treat authorization state as mutable, managed by privileged administrators, and enforceable only at the boundary of the system that manages them. When agents can spawn children, redirect delegation, or modify their own scope at runtime, the authorization state that governed a terminal action may no longer be the authorization state that was recorded at spawn time --- a gap that Gilham et al. characterize as the \emph{delegatee ambiguity problem} in multi-hop delegations \cite{hdp-draft,hdp-draft-2}.

Tamper-evident logging addresses the integrity of records but not the immutability of the delegation chain itself. Hash-chained append-only ledgers, as implemented in blockchain-based provenance systems, make it possible to detect whether records have been altered after the fact \cite{prov-standard}. However, tamper-evident logging does not by itself prevent the mutable parent-child relationships that create the ambiguity in the first place. A ledger that faithfully records a re-parenting event does not make the re-parenting legitimate --- it merely makes the illegitimate event auditable. RSP addresses this gap by making the parent pointer immutable at the protocol level, before any re-parenting event can occur, so that the delegation chain's topology is a structural constraint rather than a recorded assertion.

Recent work has formalized this problem in the agentic context. Researchers established that agent systems cannot be truly accountable without auditability as a first-class system property, identifying five dimensions of agent auditability: action recoverability, lifecycle coverage, policy checkability, responsibility attribution, and evidence integrity \cite{arxiv260405485}. Their analysis demonstrates that current agent systems widely lack the security prerequisites for auditability, and that pre-execution mediation with tamper-evident records incurs only approximately 8.3 milliseconds of median overhead --- a cost that is acceptable in high-stakes operational contexts. The SentinelAgent framework formalizes seven properties for secure delegation chains in multi-agent systems, including authority narrowing, policy preservation, forensic reconstructibility, and scope-action conformance, demonstrating that delegation chain properties can be mechanically verified via TLA+ model checking \cite{sentinelagent2025}. These results motivate the structural approach of RSP: rather than relying on post-hoc auditing to detect accountability failures, RSP makes the accountability failures structurally impossible by design.

% -------------------------------------------------------
\subsection{Fixed Versus Mutable Delegation Chains}
% -------------------------------------------------------

Delegation chains in multi-agent systems take two structurally different forms with profoundly different accountability properties. In a \emph{fixed delegation chain}, the parent-child relationship is established at spawn time and is architecturally immutable thereafter: the child's parent pointer cannot be changed, the child's authority cannot be re-delegated except through the defined spawn protocol, and the chain terminates at a predetermined anchor. In a \emph{mutable delegation chain}, the parent-child relationship is treated as runtime state that can be updated, redirected, or reassigned as conditions change.

Mutable delegation chains offer operational flexibility: a child agent can be re-parented to a different parent if the organizational structure changes, if a parent fails, or if load balancing requires reassignment. This flexibility comes at a significant accountability cost. If the parent-child relationship is mutable at runtime, then any single action by a child agent may have been authorized by one parent at one moment and a different parent at another --- and the effective authorization for a given action may not be the authorization that was recorded at spawn time. The chain of provenance becomes conditional on the state of a mutable mapping rather than fixed by an immutable record. This introduces the \emph{delegatee ambiguity problem}: after an event, it may be impossible to determine which principal effectively authorized the action, because the delegation chain has been rewritten since the action occurred.

The argument for mutability is operational necessity: in dynamic environments, task requirements change, agents fail, and organizational structures evolve. A static parent pointer becomes a liability when the parent is decommissioned. The conventional response is governance by policy: administrative documents specify conditions under which parent pointers may be redirected, and audit logs record redirection events. This approach works in low-stakes environments but fails in high-stakes ones where the cost of an unauthorized redirection may be catastrophic and irreversible. The HDP protocol identifies this precisely as the core vulnerability in conventional delegation: without cryptographic chain-of-custody, there is no verifiable link between a terminal action and its original human authorizer \cite{hdp-draft,hdp-draft-2}.

A fixed delegation chain permits no modification to the parent pointer after provisioning. The chain is established once, at spawn time, and persists for the agent's lifetime. This is a stronger constraint than access-control policy: the Fact Layer enforces immutability of the parent pointer through write protocol, rejecting any request to overwrite $\alpha(n)$ after provisioning regardless of who issued the request. The immutability is structural, not contractual: any attempt to overwrite $\alpha(n)$ is rejected at the protocol level, before execution. This contrasts with ACL-based immutability in conventional systems, where administrative compromise leads directly to immutability violation; BX3's protocol-level enforcement has no administrative override path.

The tradeoff corresponds to what the distributed systems literature terms \emph{strong consistency} versus \emph{eventual consistency} \cite{vogels2009eventual}: fixed chains correspond to strong consistency (every read of the delegation chain returns the same value, established at write time), while mutable chains correspond to eventual consistency (the value converges over time but reads during the convergence window return stale information). In low-stakes applications, eventual consistency is acceptable because the cost of divergence is low. In high-stakes agentic applications, the cost of divergence can be catastrophic and irreversible, making strong consistency the appropriate model. RSP is the structural mechanism that converts mutable delegation chains into fixed ones without sacrificing the organizational flexibility that makes multi-agent systems useful.

% -------------------------------------------------------
\subsection{Hierarchical Task Decomposition in AI}
% -------------------------------------------------------

Hierarchical task decomposition is a foundational concept in both classical AI planning and modern LLM-based agentic systems. In classical planning, a complex task is decomposed into a hierarchy of subtasks, each of which may itself be further decomposed until sub-tasks are directly executable --- a structure that enables efficient search, parallel execution, and bounded reasoning depth \cite{planner2000}. Hierarchical Task Networks (HTN) formalized this approach, requiring each decomposition to be justified by a task reduction schema that demonstrates how the composite task is equivalent to its decomposition. HTN planning's key insight is that decomposition is not merely a computational convenience but a reasoning primitive: the hierarchy of tasks encodes the causal structure of the plan.

Modern LLM-based agentic systems independently rediscovered hierarchical decomposition as a runtime pattern. Systems like AutoGPT, LangChain, and MetaGPT decompose user-provided goals into sub-tasks that are executed by child agents, which may in turn decompose their sub-tasks further. The resulting structure is a runtime-generated task tree in which depth corresponds to refinement granularity. This approach has proven effective for complex, multi-step reasoning but introduces the accountability gap that RSP addresses: each node in the task tree is a runtime entity that exercises authority, and that authority must be traceable to the human who authorized the root goal.

Satija et al. identify a related failure mode they term \emph{unnecessary depth} in task decomposition: agents may decompose a task into more steps than necessary, consuming depth without delivering additional capability \cite{zhuge2024}. This is the decomposition analog of the classical planning problem of \emph{deferred parallelizability} --- the decomposition of a task that should execute concurrently into a sequential chain consuming more depth without delivering more capability. The Purpose Layer's spawn necessity check forces the parent to articulate why spawning is necessary, making artificial fragmentation auditable and rejectable.

Recursive Spawning adopts the hierarchical decomposition paradigm from both traditions but adds the critical accountability layer absent from both: each node in the decomposition tree is a first-class accountability entity with an immutable parent pointer, a defined termination condition, and a cryptographically verifiable chain back to a human anchor. Decomposition is not merely a planning technique; it is an authorization event that creates a new accountable entity. The decomposition tree is therefore not just a planning artifact but an organizational chart of delegated authority.

% -------------------------------------------------------
\subsection{The BX3 Framework as Substrate for Immutable Parent Chains}
% -------------------------------------------------------

The BX3 Framework provides the architectural substrate within which Recursive Spawning operates as a first-class, accountable delegation primitive \cite{bx3framework2026}. The BX3 Framework organizes any autonomous node into three immutable functional layers --- the \textit{Purpose Layer}, the \textit{Bounds Engine}, and the \textit{Fact Layer} --- each defined by the properties it must maintain rather than by the type of actor occupying it. The Purpose Layer carries intent, judgment, and final human accountability; the Bounds Engine carries constrained execution within a defined operating range; the Fact Layer carries deterministic enforcement and forensic auditability. Together they convert a spawn chain from a flexible delegation heuristic into a verifiable accountability infrastructure.

Three BX3 layers serve as the backbone of every RSP spawn chain. The \textit{Purpose Layer} declares and anchors intent. It defines the maximum permissible chain depth as part of the spawn authorization policy and records the identity of the human principal who authorized the root spawn event, propagating this anchor ($h(n)$) to all descendants. At every spawn event, the Purpose Layer requires a scope refinement justification and a spawn necessity declaration, ensuring that each new child represents a genuine refinement of its parent's authorized scope rather than an independent expansion of authority. The \textit{Bounds Engine} enforces structural and computational constraints at every spawn event. It operates along two axes: depth management, which enforces the maximum chain depth $D_{\max}$ defined by the Purpose Layer, and scope management, which enforces the subset relationship constraint that each child's operating scope $R(n_{i+1})$ must be a proper subset of its parent's $R(n_i)$. A spawn that would violate either constraint is refused regardless of the parent agent's intentions. The \textit{Fact Layer} is the append-only, tamper-evident forensic ledger that records every meaningful event in a spawn chain --- every spawn event, state transition, Purpose Layer directive, and exception condition. Each ledger entry is hash-chained to its predecessor, making the ledger tamper-evident: any insertion, deletion, or modification breaks the hash chain and is detectable by any observer with read access. The Fact Layer is also the layer that produces the immutable parent pointer: $\alpha(n)$ is established once, at node provisioning, through a Fact Layer write operation that is atomic and immediately hash-chained to the preceding ledger entry. After this write, no actor --- including the Purpose Layer that originated it --- can modify $\alpha(n)$ because the Fact Layer's write protocol rejects any modification to authorization records after they are established.

The upstream accountability guarantee of the BX3 Framework states that when any node encounters a state it cannot resolve within its provisioned bounds, accountability escalates recursively upward through the system hierarchy, bypassing all machine actors, until it reaches a human anchor \cite{bx3framework2026}. This guarantee is structural, not policy-driven: the Fact Layer creates the immutable parent pointer at spawn time, issues the spawn warrant, and maintains the forensic ledger that records every spawn event. The structural necessity of the three-layer mapping is established formally in Section~\ref{sec:rs-integration}. Removing any layer causes the corresponding guarantee to fail. Without the Purpose Layer, there is no authoritative human anchor and no source of spawn authorization policy. Without the Bounds Engine, depth limits become unenforceable conventions subject to override by operational urgency. Without the Fact Layer, immutability of the parent pointer and tamper-evidence of the ledger are promises rather than constraints. The three-layer architecture is not a design pattern --- it is the minimal structure that makes RSP's correctness properties unconditionally guaranteed rather than statistically unlikely.

% -------------------------------------------------------
\subsection{Synthesis}
% -------------------------------------------------------

The five research traditions reviewed above converge on a common diagnosis: conventional agentic architectures lack the structural mechanisms needed to guarantee that every agent in a spawn chain operates within a delegation chain that is simultaneously immutable, auditable, scope-refined, depth-bounded, and verified to terminate at a human principal. HDP, ALP, PROV-AGENT, UALM, and SentinelAgent each contribute a necessary component of the solution --- cryptographic delegation tokens, explicit lifecycle state machines, W3C PROV provenance modeling, five-layer fleet governance, and formal delegation chain calculi --- but none of them individually provides the structural immutability of the parent pointer combined with the depth bound, scope subset enforcement, and tamper-evident forensic ledger in a single unified protocol. RSP is the structural mechanism that integrates these components within the BX3 Framework's three-layer architecture, producing an accountable spawn chain that is structurally impossible to corrupt rather than merely audited after the fact.
